% The pricing procedure for the works on the website.
% Build: latexmk -pdf pricing-procedure.tex   (pdfLaTeX; build files are git-ignored)
% The numbers below are produced by src/lib/price-formula.ts and checked by
% tests/unit/price-formula.test.ts: change the three together.
\documentclass[11pt,a4paper]{article}

\usepackage[T1]{fontenc}
\usepackage{lmodern}
\usepackage[margin=2.4cm]{geometry}
\usepackage{microtype}
\usepackage{amsmath}
\usepackage{booktabs}
\usepackage{array}
\usepackage[table]{xcolor}
\usepackage{enumitem}
\usepackage[most]{tcolorbox}
\usepackage[hidelinks]{hyperref}

\definecolor{lake}{HTML}{2E7D8C}
\newcommand{\nis}[1]{NIS\,#1}
\newcommand{\code}[1]{\nolinkurl{#1}}  % breaks at / and .
\newcommand{\Pref}{P_{\mathrm{ref}}}
\newcommand{\Aref}{A_{\mathrm{ref}}}
\setlist{itemsep=2pt,topsep=4pt}

\title{\textbf{Pricing procedure}\\[4pt]\large How every price of Geula Levy's works on the website
  is set}
\author{Geula Levy website}
\date{October 2026}

\begin{document}
\maketitle

\begin{tcolorbox}[colback=lake!6,colframe=lake,title=\textbf{The rule},fonttitle=\sffamily]
Every work is priced by one formula, so that works of the same size and material always cost the
same, and larger works cost more but less than in proportion to their area:
\[
  P \;=\; \operatorname{round}_{100}\!\Bigl(\Pref \cdot m_{\text{medium}} \cdot m_{\text{surface}}
  \cdot \Bigl(\frac{A}{\Aref}\Bigr)^{\beta}\Bigr) \;+\; F_{\text{framed}},
  \qquad P_{\text{USD}} \;=\; \operatorname{round}_{10}\!\Bigl(\frac{P}{3.7}\Bigr)
\]
with $\Pref = \nis{4{,}800}$ (an $80\times80$\,cm acrylic on canvas), $\Aref = 6{,}400$\,cm$^2$,
$\beta = 0.6$, and $A$ the painted area in cm$^2$.
\end{tcolorbox}

\section{Where the formula lives}
\begin{itemize}
  \item \textbf{Code:} \code{src/lib/price-formula.ts} holds the parameters (\code{PRICING}), the
    material factors and the two functions \code{formulaPriceIls} and \code{formulaPriceUsd}.
  \item \textbf{Seed:} \texttt{npm run db:reset} (the \code{demo} seed, the painter's catalog) prices
    every work with the formula (\code{asCatalogWork} in \code{scripts/lib/catalog-manifest.ts}).
    The catalog manifest stores the same prices for reference.
  \item \textbf{Tests:} \code{tests/unit/price-formula.test.ts} checks the formula against every
    number in this document and checks that the manifest matches it.
  \item \textbf{Exception:} the 16 test works of the \code{fixtures} seed keep their own test
    prices. They are used only by the automated tests, never on the real site.
\end{itemize}

\section{Parameters}
\begin{center}
\small
\rowcolors{2}{gray!6}{white}
\begin{tabular}{@{}l>{\raggedright\arraybackslash}p{5.6cm}r>{\ttfamily}l@{}}
  \toprule
  Symbol & Meaning & Value & \normalfont In \code{PRICING} \\
  \midrule
  $\Pref$ & price of the reference work, an $80\times80$\,cm acrylic on canvas & \nis{4{,}800} & referenceIls \\
  $\Aref$ & area of the reference work & 6{,}400\,cm$^2$ & referenceAreaCm2 \\
  $\beta$ & size exponent: price grows with area$^{\beta}$ & 0.6 & sizeExponent \\
  & rounding of the shekel price & \nis{100} & roundIls \\
  $F_{\text{framed}}$ & added for a work sold framed (the frame's cost) & \nis{0} for now & framedAddIls \\
  & reference rate for the dollar price & 3.7\,NIS/\$ & ilsPerUsd \\
  & rounding of the dollar price & \$10 & roundUsd \\
  \bottomrule
\end{tabular}
\end{center}

\section{Material factors}
The material factor is the medium factor times the surface factor. Acrylic on canvas, 148 of her
151 works, is the base.
\begin{center}
\begin{tabular}{@{}llr@{}}
  \toprule
  Factor & Material & Value \\
  \midrule
  $m_{\text{medium}}$ & oil & 1.15 \\
                      & every other medium (acrylic, watercolour, tempera, mixed media, \ldots) & 1.00 \\
  \midrule
  $m_{\text{surface}}$ & canvas, linen & 1.00 \\
                       & wood panel, board, cardboard & 0.95 \\
                       & paper & 0.65 \\
  \bottomrule
\end{tabular}
\end{center}
For an $80\times80$ work: acrylic on canvas \nis{4{,}800}, oil on canvas \nis{5{,}500}, acrylic on
board \nis{4{,}600}, watercolour on paper \nis{3{,}100}.

\section{Pricing a work, step by step}
\begin{enumerate}
  \item Measure the painted surface: height $H$ and width $W$ in centimetres, without the frame.
    Orientation does not matter.
  \item Area: $A = H \times W$ (cm$^2$).
  \item Material factor: $m = m_{\text{medium}} \times m_{\text{surface}}$ (Section~3).
  \item Raw price: $4{,}800 \times m \times (A / 6{,}400)^{0.6}$.
  \item Round to the nearest \nis{100}.
  \item If the work is sold framed, add $F_{\text{framed}}$ (currently \nis{0}: the frame is part of
    the work).
  \item Dollar price: divide by 3.7 and round to the nearest \$10.
  \item In the admin: \emph{Artworks} $\rightarrow$ the work $\rightarrow$ \emph{Price \& sale}.
    Enter \emph{Price in ILS} and \emph{Price in USD}, and tick \emph{I confirm the price change}.
\end{enumerate}

\paragraph{Example 1: \emph{By the Lake}, acrylic on canvas, $60\times80$.}
$A = 4{,}800$; $(4{,}800/6{,}400)^{0.6} = 0.75^{0.6} = 0.8415$; $4{,}800 \times 0.8415 = 4{,}039$
$\Rightarrow$ \nis{4{,}000}; $4{,}000/3.7 = 1{,}081$ $\Rightarrow$ \$1{,}080.

\paragraph{Example 2: an oil on canvas, $60\times60$.}
$A = 3{,}600$; $(3{,}600/6{,}400)^{0.6} = 0.5625^{0.6} = 0.7080$;
$4{,}800 \times 1.15 \times 0.7080 = 3{,}908$ $\Rightarrow$ \nis{3{,}900};
$3{,}900/3.7 = 1{,}054$ $\Rightarrow$ \$1{,}050.

\section{Price list}
Prices for common sizes, unframed (sizes in cm, area in cm$^2$):
\begin{center}
\rowcolors{2}{gray!6}{white}
\begin{tabular}{@{}lrrrrr@{}}
  \toprule
  & & \multicolumn{2}{c}{Acrylic on canvas} & \multicolumn{2}{c}{Oil on canvas} \\
  \cmidrule(lr){3-4}\cmidrule(l){5-6}
  Size & Area & NIS & USD & NIS & USD \\
  \midrule
  40\,$\times$\,60 & 2{,}400 & 2{,}700 & 730 & 3{,}100 & 840 \\
  60\,$\times$\,60 & 3{,}600 & 3{,}400 & 920 & 3{,}900 & 1{,}050 \\
  60\,$\times$\,80 & 4{,}800 & 4{,}000 & 1{,}080 & 4{,}600 & 1{,}240 \\
  70\,$\times$\,100 & 7{,}000 & 5{,}100 & 1{,}380 & 5{,}800 & 1{,}570 \\
  80\,$\times$\,80 & 6{,}400 & 4{,}800 & 1{,}300 & 5{,}500 & 1{,}490 \\
  80\,$\times$\,100 & 8{,}000 & 5{,}500 & 1{,}490 & 6{,}300 & 1{,}700 \\
  90\,$\times$\,90 & 8{,}100 & 5{,}500 & 1{,}490 & 6{,}400 & 1{,}730 \\
  80\,$\times$\,120 & 9{,}600 & 6{,}100 & 1{,}650 & 7{,}000 & 1{,}890 \\
  100\,$\times$\,100 & 10{,}000 & 6{,}300 & 1{,}700 & 7{,}200 & 1{,}950 \\
  100\,$\times$\,120 & 12{,}000 & 7{,}000 & 1{,}890 & 8{,}000 & 2{,}160 \\
  100\,$\times$\,150 & 15{,}000 & 8{,}000 & 2{,}160 & 9{,}200 & 2{,}490 \\
  120\,$\times$\,145 & 17{,}400 & 8{,}700 & 2{,}350 & 10{,}100 & 2{,}730 \\
  \bottomrule
\end{tabular}
\end{center}
International parcels are capped at \$2{,}500 declared value with DHL (\emph{Settings}
$\rightarrow$ \emph{Shipping}). A work priced above the cap, such as the largest oil, is offered
abroad by quote only.

\section{Rules}
\begin{enumerate}
  \item \textbf{One price per size and material.} Subject and series do not change the price. An
    exceptional work may be set up to 15\,\% above or below the formula, rarely, and for a reason a
    buyer would understand (unusual complexity, a key work of a series).
  \item \textbf{One price everywhere.} The website, the studio and any future gallery show the same
    price. A gallery takes its commission from that price; it does not add to it.
  \item \textbf{Prices only go up.} When 70--80\,\% of a body of work has sold, raise $\Pref$ by
    10--15\,\%. Every price moves by the same percentage, so the list stays consistent. Never lower
    prices.
  \item \textbf{VAT.} The admin's shekel price is the price the buyer pays. As an \emph{osek patur}
    she charges no VAT. If she becomes an \emph{osek murshe}, that price includes VAT (18\,\%):
    either keep the prices and absorb it, or raise $\Pref$ by 18\,\%.
  \item \textbf{Dollar prices} follow the shekel prices at the reference rate. If the rate moves a
    lot, update \code{ilsPerUsd} and the dollar prices together.
  \item \textbf{Shipping and insurance} are charged separately at checkout and are not part of the
    price.
\end{enumerate}

\section{Changing the formula}
\begin{enumerate}
  \item Edit \code{PRICING} or a factor in \code{src/lib/price-formula.ts}. To raise all prices by
    12\,\%, for example, set \code{referenceIls} to 5{,}400 (4{,}800 $\times$ 1.12 = 5{,}376,
    rounded).
  \item Recompute the manifest prices and the examples in this document. Run
    \code{tests/unit/price-formula.test.ts}; it fails until both match the new formula.
  \item A local database picks up the new prices with \texttt{npm run db:reset}. A live database is
    not repriced automatically: change the prices in the admin, or have a developer run a one-off
    update that applies the formula to the works that have no individual exception.
\end{enumerate}

\section{Why this formula}
Art-market research that models auction prices from a work's characteristics finds that price
rises with size, but less than in proportion. A 10\,\% larger canvas sells for roughly 3--7\,\%
more \cite{masterworks,renneboog}, and about 6\,\% for contemporary painters (an elasticity of
0.59 in a 2026 study of 18{,}821 works) \cite{orstavik}.
\begin{itemize}
  \item \textbf{The exponent 0.6} follows that research. Pricing by area (``per square inch'') has an
    exponent of~1 and overprices large works. Pricing by width plus height (``linear inch'') has
    0.5 for a given shape and is recommended for its gentler steps \cite{foundmyself}.
  \item \textbf{Works on paper} sell 20--50\,\% below canvas \cite{milan,masterworks}, hence the
    0.65 factor.
  \item \textbf{Oil} sells above acrylic; a common set of multipliers gives oil 1.30 and acrylic 1.10,
    a ratio of about 1.18 \cite{priceyourpainting}, rounded down here to 1.15.
  \item \textbf{The reference price} \nis{4{,}800} for $80\times80$ (about \nis{0.75} per cm$^2$) is
    at the upper end of published direct-sale rates for artists without a gallery or auction
    record. Examples: \$1 per in$^2$ direct, about \nis{3{,}700} for this size \cite{milan};
    \pounds1.00--1.25 per in$^2$, about \nis{4{,}700}--\nis{5{,}800} \cite{priceyourpainting}.
    It is below gallery and established-artist levels.
\end{itemize}

{\raggedright
\begin{thebibliography}{9}
\bibitem{masterworks} Masterworks Research, \emph{Hedonic regression in art pricing: how
  characteristics drive value}, June 2026.
  \url{https://www.masterworks.com/academy/posts/hedonic-regression-in-art-pricing-how-characteristics-drive-value}
\bibitem{renneboog} L.~Renneboog and C.~Spaenjers, Buying beauty: on prices and returns in the art
  market, \emph{Management Science} 59(1):36--53, 2013.
\bibitem{orstavik} O.~H.~S.~{\O}rstavik, \emph{Size Matters: How Physical Scale Is (or Is Not)
  Rewarded in a Global, Ultra-Contemporary, Online Art Market}, MA thesis, Sotheby's Institute of
  Art, 2026. \url{https://digitalcommons.sia.edu/ma_art_bus/22}
\bibitem{foundmyself} FoundMyself, \emph{How to price your art}.
  \url{https://www.foundmyself.com/blog/how-to-price-your-art/}
\bibitem{milan} Milan Art Institute, \emph{Sell \& Price Your Art}.
  \url{https://blog.milanartinstitute.com/how-to-price-paintings/}
\bibitem{priceyourpainting} Price Your Painting, \emph{How to value a painting: the complete guide
  for artists}. \url{https://priceyourpainting.substack.com/p/how-to-value-a-painting-the-complete}
\end{thebibliography}
}

\end{document}
